\subsection*{IV 20.5: Théorèmes du type Seidenberg (cote 25, sites 22--25, with the restored paragraph from site 21)}
The opening instruction restores the preceding paragraph about openness, whose margin directs it to no.\ 5.1. The section then uses the incidence divisor and its defining section $\phi$, expresses the inverse-image module as a cokernel, and defines $F$-regularity. Proposition 5.3 treats the regularity locus $U$ and relative flatness; corollary 5.4 gives the corresponding locus $V$ in the dual projective space. The pen loop on site 23 moves the finite-presentation and flatness hypothesis into part a), and the handwritten additions to 5.4 distinguish its constructibility and openness assertions. Remark 5.5 discusses associated cycles and supports; theorem 5.6 and corollary 5.7 invoke the lists in 12.1.1 and 12.1.4. The text expressly uses 11.3 for the flatness argument, mentions EGA III without a completed number in the locally noetherian case, and refers to paragraph 3 and nos.\ 2.3 and 2.6 for the associated-cycle argument. Site 24 asks that the final exposition replace references to lists in another fascicule by an explicit list and identify Seidenberg's geometrically normal case. The crossed commentary on site 25 explains the proposed division of the preceding material and compares the route through 4.6 with the ``délicat 15.2.2''. It also states that the author has not read Seidenberg's paper. The checked text retains mathematical slips visible in the formulation: site 22 speaks of an isomorphism of $F$ with $\underline{O}_X$ after discussing an invertible $L$ and later writes $\mathrm{id}_F\otimes\phi$ before the displayed source $F\otimes L^{-1}$; site 23 prints $s\in\check P$ in 5.3 b); site 24 concludes with flatness relative to $S$ in 5.6. The cancelled handwritten paragraph on site 23 reads ``$k$ un corps'', correcting the draft transcription's ``$A$ complet''; its comparison sign remains uncertain. The cancelled Remords paragraph on site27 returns to the same proposed reorganization: it separates openness proved through a valuative criterion from the additional flatness conclusion using15.2.2, and proposes a new4.8. Sources: \href{https://grothendieck.umontpellier.fr/web/viewer.html?file=../25.pdf\#page=22}{site 21}, \href{https://grothendieck.umontpellier.fr/web/viewer.html?file=../25.pdf\#page=23}{sites 22--25}; published EGA IV references explicitly written here include 11.3, 12.1.1, 12.1.4 and 15.2.2.
